\chapter{PBL Problem Map}

\ifdefined\colA\else\newlength{\colA}\fi
\ifdefined\colB\else\newlength{\colB}\fi
\ifdefined\colC\else\newlength{\colC}\fi
\setlength{\colA}{\dimexpr 0.13\textwidth - \tabcolsep\relax}
\setlength{\colB}{\dimexpr 0.53\textwidth - \tabcolsep\relax}
\setlength{\colC}{\dimexpr 0.34\textwidth - 2\tabcolsep\relax}

% Reusable command for a single-row table
% #1 = col1, #2 = col2, #3 = col3, #4 = extra rule (pass \midrule or nothing)
\newcommand{\pblrow}[4]{%
  \noindent{\small\begin{tabular}{@{}
    >{\raggedright\arraybackslash}p{\colA}
    >{\raggedright\arraybackslash}p{\colB}
    >{\raggedright\arraybackslash}p{\colC}@{}}
  #4
  \parbox[t]{\colA}{\small #1\strut} &
  \parbox[t]{\colB}{\small #2\strut} &
  \parbox[t]{\colC}{\small #3\strut} \\
  \bottomrule
  \end{tabular}}%
  \par\medskip
}

% ── Header row ───────────────────────────────────────────────────────────────
\noindent{\small\begin{tabular}{@{}
  >{\raggedright\arraybackslash}p{\colA}
  >{\raggedright\arraybackslash}p{\colB}
  >{\raggedright\arraybackslash}p{\colC}@{}}
\toprule
\textbf{TYPE OF QUESTION} & \textbf{EXAMPLES OF QUESTIONS} & \textbf{SHORT ANSWER} \\
\bottomrule
\end{tabular}}
\par\medskip

% ── Row a) ───────────────────────────────────────────────────────────────────
\pblrow{a)}{Present the problem situation in 5--7 sentences.}{%
  Children today face unprecedented challenges with attention span and sustained
  focus, making traditional programming education---which often involves lengthy
  tutorials, abstract concepts, and delayed gratification---increasingly
  ineffective. Gamification has emerged as a powerful educational strategy that
  leverages game mechanics, instant feedback, and intrinsic motivation to
  maintain engagement and facilitate learning. Minecraft represents an ideal
  platform for gamified programming education, as children are already deeply
  invested in the game and motivated to enhance their gameplay experience.
  However, existing Minecraft programming solutions are too complex for young
  learners and fail to incorporate gamified learning principles that could
  sustain their attention. There is a critical need for a domain-specific
  language that transforms programming education into an engaging, game-like
  experience within Minecraft, where children see immediate results from their
  code through scripted actions executed in the game.%
}{\toprule}

% ── Row b) ───────────────────────────────────────────────────────────────────
\pblrow{b)}{Formulate the problem -- 1 statement.}{%
  Children's limited attention spans prevent them from learning programming
  through traditional methods, which often lack immediate visual feedback and
  engaging concepts.%
}{\toprule}

% ── Who ──────────────────────────────────────────────────────────────────────
\pblrow{Who\ldots?}{%
  1. Who has a problem?\par
  2. Who is involved in the problem situation?\par
  3. Who suffers from the problem?\par
  4. Who should take part in solving the problem?%
}{%
  1. Children.\par
  2. Children, parents, professionals of the educational domain,
     authorities concerned with education.\par
  3. Children.\par
  4. Authorities concerned with education, DSL designers and developers,
     teachers.%
}{\toprule}

% ── What / Which ─────────────────────────────────────────────────────────────
\pblrow{What\ldots?\par Which\ldots?}{%
  5. What happened?\par
  6. What is the problem?\par
  7. What led you to identify the problem?\par
  8. What are the causes that affect those involved?\par
  9. What are the needs of those targeted?\par
  10. What resources are necessary to solve the problem?%
}{%
  5. Short-form media has become a main source of information.\par
  6. Incompatibility between traditional programming education and the
     attention patterns of modern children.\par
  7. Research showing declining attention spans in children.\par
  8. Fast-paced media consumption; limited availability of age-appropriate
     engaging programming tools.\par
  9. Immediate visual feedback, game-based learning, progressive difficulty,
     creative freedom.\par
  10. DSL developers, educational consultants, UI/UX designers, Minecraft
      experts, child testers, teacher advisors.%
}{\toprule}

% ── Where ────────────────────────────────────────────────────────────────────
\pblrow{Where\ldots?}{%
  11. Where did the problem arise?%
}{%
  11. At the intersection of children's educational environments (homes,
      schools, after-school programs) and their recreational gaming spaces.%
}{\toprule}

% ── When ─────────────────────────────────────────────────────────────────────
\pblrow{When\ldots?}{%
  12. When did the problem arise?\par
  13. When was the problem identified?\par
  14. When is it planned to be resolved?%
}{%
  12. Progressively over the past decade (2015--2025), as exposure to
      fast-paced digital media increased.\par
  13. During the initial stages of the PBL project (2025--2026 academic year).\par
  14. By the end of the current semester.%
}{\toprule}

% ── Why ──────────────────────────────────────────────────────────────────────
\pblrow{Why\ldots?}{%
  15. Why is there a need to solve the problem?%
}{%
  15. Programming and computational thinking are essential 21st-century skills.
      Without intervention, an entire generation may miss foundational digital
      literacy skills necessary for future academic and career success.%
}{\toprule}

% ── How ──────────────────────────────────────────────────────────────────────
\pblrow{How\ldots?}{%
  16. How can the problem be solved?\par
  17. How will the elimination of the problem be determined?%
}{%
  16. Gamification of the learning process.\par
  17. By the number of users of the solution and the number of educational
      systems that have adopted it.%
}{\toprule}

% ── How much ─────────────────────────────────────────────────────────────────
\pblrow{How much\ldots?}{%
  18. How big is the problem?\par
  19. How many people are affected?\par
  20. How long will it take to resolve?\par
  21. How many parties need to be involved?%
}{%
  18. Global~--- affects generations of children with excessive exposure to
      digital technologies.\par
  19. Estimated at over 140 million children.\par
  20. Implementation could potentially be completed in 6~months.\par
  21. Minimum 3 parties.%
}{\toprule}

% ── Other ────────────────────────────────────────────────────────────────────
\pblrow{Other (optional)}{%
  22. How do you see solving this problem?\par
  23. Conclusions and reflections.%
}{%
  22. The product should contain simplified programming syntax, educator
      resources, and use-case tutorials suited for children. Given Minecraft's
      popularity, success potential is high.\par
  23. The Minecraft Bot Maker DSL addresses a genuine and growing educational
      challenge; if done right, it could become a fundamental learning tool.%
}{\toprule}

\clearpage